\chapter{Background}
\label{cha:background}

This chapter provides the technical background needed to motivate the system design in Chapter~\ref{cha:design}. Section~\ref{sec:bg_kvcache} covers the latency bottleneck caused by redundant prefill computation in agentic workflows and presents CacheBlend as the state-of-the-art approach to position-independent KV cache reuse. Section~\ref{sec:bg_context} covers context retrieval as the complementary frontend concern and presents RepoGraph as the concrete retrieval module used in our integrated system. Section~\ref{sec:bg_motivation} explains why an integrated system is needed to evaluate the backend enhancements under realistic agentic workloads.

\section{Position-Independent KV Cache Reuse}
\label{sec:bg_kvcache}

A primary source of latency and resource inefficiency in agentic LLM serving is redundant computation across requests. Many agentic frameworks send a large number of sequential requests to the LLM during a single task. The prompt in each request typically includes bulk context information and template tokens. When no cache reuse is in place, the model re-processes these tokens from scratch every time. For instance, an agent might retrieve a large snippet of code, feed it into the model to propose a fix, and then in a subsequent step feed much of the same snippet again while testing or refining the fix. Each of these calls recomputes attention over the repeated context tokens. This redundant computation is a shared inefficiency across many agentic workflows.

% [RESOLVED] \rashmi{Add ref}
At the core of this inefficiency is the key-value (KV) cache in transformer models, which stores intermediate attention states of processed tokens so that each new token need not recompute attention over the entire sequence. The cache normally resets between separate API calls. To reuse computation across calls, several optimized LLM serving engines have introduced caching and scheduling strategies. vLLM \cite{kwon2023vllm} improves throughput by batched execution and paged memory management of the KV cache, treating it as pageable memory. It can serve multiple generation requests in parallel and merge requests with identical prefixes so that the shared prefix is computed only once. Even vLLM, however, only reuses prefix computations when those prefixes exactly align from the start of the prompt \cite{kwon2023vllm, zheng2024sglangefficientexecutionstructured}. In an agent scenario, calls often share partial contexts that appear after some unique tokens, leading to low cache hit rates. SGLang \cite{zheng2024sglangefficientexecutionstructured} introduces RadixAttention, storing caches in a radix tree keyed by prefixes and scheduling requests to maximize cache hits, though it too relies on exact prefix matching. Another approach, Prompt Cache \cite{gim2024promptcachemodularattention}, allows explicitly marking reusable prompt segments and precomputing their KV states. While this yields substantial latency reductions on long-document QA tasks, it requires manual template setup and cannot adapt if templates change.

Cache optimization for LLM inference has recently taken a leap forward with CacheBlend \cite{yao2025cacheblendfastlargelanguage}. CacheBlend's key innovation is enabling \emph{position-independent} reuse of KV caches. It treats the prompt as composed of chunks, checks for each chunk whether it has stored KV state, and when reassembling chunks selectively recomputes a small subset of tokens. The recomputation target is the set of HKVD (High-KV-Deviation) tokens, which are tokens whose cached KV state would otherwise deviate too far from what the model would have computed if the prompt were processed from scratch. Recomputing these tokens restores the effect of cross-attention on the forward-attention matrix. The observation that HKVD tokens tend to remain deviated through different layers also makes it possible to select and recompute them efficiently. On the multi-round QA workloads used in the original paper, this blending enables much higher cache hit rates and yields a reported 2--3$\times$ reduction in time-to-first-token and 3--5$\times$ throughput improvement without degrading output quality. CacheBlend was demonstrated effective in multi-round QA scenarios through datasets such as 2WikiMQA and Musique \cite{ho2020constructingmultihopqadataset, trivedi2022musiquemultihopquestionssinglehop}.

% [RESOLVED] \rashmi{Too much detail. This looks like a bug fix. If so, need to deemphaize this.}
CacheBlend's reference implementation, LMCache, leaves several aspects of agentic workloads unaddressed. First, the mechanism that corrects positional encodings when a chunk is reused at an offset different from where it was originally cached is incomplete, so position-shifted reuse does not work in practice. As a concrete consequence, if a code snippet appears at positions $[120,\,180]$ in one request and at positions $[400,\,460]$ in the next, the cached KV state from the first request cannot be applied directly to the second, because the rotary embeddings were computed for the original positions and need to be corrected for the new ones. Second, the storage backend allocates one memory object per cache chunk, acquiring the in-memory CPU storage lock on each allocation. For a prompt with $N$ chunks this serializes $N$ lock acquisitions, which becomes a bottleneck under high concurrency. Third, the eviction loop on a full cache retries indefinitely, which can deadlock when all entries are pinned by concurrent lookups on other threads. Finally, the KV store runs synchronously inside the model execution step, so a request pays an offload cost proportional to its prompt length inside its own time-to-first-token, even when it reads nothing from the cache. Chapter~\ref{cha:design} presents the modifications introduced to address the first three. The fourth is identified by our evaluation in Chapter~\ref{cha:eval}, which quantifies its cost and shows that it dominates the aggregate result.

\section{Repository-Aware Context Retrieval}
\label{sec:bg_context}

While cache reuse on the backend addresses latency, the quality of the output the agent produces depends on the context that is included in the prompt. This is a separate concern from the backend, but a complete end-to-end system requires some method for assembling per-request prompts that contain the right code context for a given task. We briefly review the state of context retrieval here because the frontend used in our evaluation builds on this prior work.

% [RESOLVED] \rashmi{ This may not be true anymore. Just write this as how the field evolved rather than saying it is still a problem. That is, in past tense.}
Early code LLM agents demonstrated strong capabilities on isolated coding tasks, but they worked from incomplete or irrelevant context when the task spanned a large repository. Frameworks such as SWE-Agent \cite{yang2024sweagent}, AutoCodeRover \cite{zhang2024autocoderover}, and Agentless \cite{xia2024agentlessdemystifyingllmbasedsoftware} fed the model only limited snippets of a codebase or generic summaries. These agents reached the rest of the codebase by searching keywords, files, and directories, together with a scrolling mechanism for file viewing. SWE-Agent, for example, exposed three commands: \verb|goto|, \verb|scroll_up|, and \verb|scroll_down|, which let the agent move a constant-sized context window to any line of the target file \cite{yang2024sweagent}. This linear view of the repository left the agent with false assumptions on code structure and implementation details, leading to unstable patch quality. Providing an agent with only a function implementation but not the definitions of its referenced symbols omitted crucial information and hurt the quality of patches. Supplying too much code instead overwhelmed the model with irrelevant details and interfered with its reasoning \cite{zhang2024hierarchicalcontextpruningoptimizing, zhang2024inftybenchextendinglongcontext}. This tension motivated the repository-aware retrieval work that followed.

% [RESOLVED] \rashmi{ What is "tree-sitter". Don't go to the level of describing the codebase in the thesis. Explain conceptually.}
Several lines of work have addressed this tension. SWE-Search \cite{antoniades2025swesearchenhancingsoftwareagents} integrates Monte Carlo Tree Search to explore the solution space, and CodePlan \cite{bairi2023codeplanrepositorylevelcodingusing} tracks incremental code changes with dependency analysis. CoSIL \cite{jiang2025cosilsoftwareissuelocalization} dynamically constructs a module-level call graph during the agent's search, expanding the graph by following function calls and pruning irrelevant branches. RepoGraph \cite{ouyang2025repographenhancingaisoftware} takes a complementary approach by pre-building a complete repository graph before the agent starts working. It constructs a repository-level code graph that captures relationships between code elements across different files: each class, method, or function is a node, and the edges connect definitions to their usage or dependency. During a setup phase, RepoGraph parses the source files to identify where each symbol is defined and where it is referenced, links each definition to its uses, and stores the resulting graph for fast querying. When the agent needs context for a particular function or error, it queries this graph to retrieve all related definitions and usage chains. RepoGraph has been shown to consistently boost the resolve rate of Agentless and SWE-Agent when paired with capable backends such as GPT-4o or Claude-3.5-Sonnet \cite{ouyang2025repographenhancingaisoftware}.

We use RepoGraph as the concrete context retrieval module in our integrated system. The choice is pragmatic rather than essential to the contribution: RepoGraph produces structured, chunk-friendly prompts that expose natural cache boundaries, and its open-source implementation is straightforward to integrate. The adapter design in Section~\ref{sec:design_arch_backend} keeps the frontend interface narrow, so any alternative retrieval module that produces prompts with explicit chunk boundaries (for example a more recent agentic-search-based system) can be substituted without affecting the backend.

\section{Motivation for Integration}
\label{sec:bg_motivation}

The original CacheBlend evaluations used multi-round question-answering datasets in which prompts share structured text chunks across turns. While these workloads validate the position-independent reuse mechanism in principle, they leave open the question of how the system performs under code-agent workloads, where prompts can consist of long code context, repository structure information, and reordered context blocks across requests. To evaluate this in a realistic setting, we wrap LMCache (with the backend enhancements introduced in Chapter~\ref{cha:design}) in an end-to-end agentic pipeline that uses RepoGraph and Agentless as the frontend to generate realistic prompts. The backend enhancements are the central technical contribution; the integration with the frontend serves as the evaluation harness that exposes the cache engine to representative agentic traffic.

A second motivation for the integrated system is a quality control point. Position-independent cache reuse trades a small amount of computed-from-scratch fidelity for latency savings when the cache hits, and it is important to confirm that the resulting outputs remain usable in practice. Our frontend evaluation in Section~\ref{sec:ch1_eval} addresses this question directly by comparing the resolve rate of the integrated system under different blending granularities against a no-blending baseline.

% [RESOLVED] \rashmi{It is unclear what "Cached Context" is at this point, despite the terse description given here. }
The mechanism that ties the backend cache engine to the frontend prompts is what we refer to as \emph{Cached Context}. Rather than treating a prompt as a single unit that must line up end to end with an earlier one, the frontend splits each prompt it produces into separate cache units at the boundaries of individual context objects: one code snippet, one repository structure block, or one template section. The backend then matches and reuses each unit on its own, so reuse no longer depends on whole prompts agreeing. Cached Context exploits the modular and repetitive structure of code-agent context, which consists mostly of repository structure information and function or class definition blocks, alongside prompt templates whose tokens recur on every request. Section~\ref{sec:design_cc} formalizes this concept, and Chapter~\ref{cha:design} presents the full system design.
